\chapter{研究方法 Methods}
\label{ch:method}

本章示範學位論文常用的三種元素：數學公式、圖、表。

\section{數學公式 Equations}
行內公式如 $E = mc^2$；編號的展示公式如式\eqref{eq:demo}：
\begin{equation}
  \label{eq:demo}
  \int_{a}^{b} f(x)\,\mathrm{d}x = F(b) - F(a).
\end{equation}

\section{圖 Figures}
以 \texttt{\textbackslash includegraphics} 置入圖檔（存於 \texttt{figures/}）。此處以
一個佔位方框代替真實圖檔，讓範本不需附圖即可編譯；實務上請改用註解中的寫法。
\begin{figure}[htbp]
  \centering
  % 真實用法：\includegraphics[width=0.7\textwidth]{figures/your-figure.pdf}
  \fbox{\begin{minipage}[c][4cm][c]{0.7\textwidth}\centering（範例圖檔佔位 placeholder figure）\end{minipage}}
  \caption{圖題範例。若圖擷取自參考文獻，須於圖題標註來源（規範八）。}
  \label{fig:demo}
\end{figure}

\section{表 Tables}
表格建議使用 \texttt{booktabs} 的三線表樣式，見表\ref{tab:demo}。
\begin{table}[htbp]
  \centering
  \caption{表題範例。若表擷取自參考文獻，須於表中標註來源（規範九）。}
  \label{tab:demo}
  \begin{tabular}{lcc}
    \toprule
    項目 Item & 數值 Value & 單位 Unit \\
    \midrule
    樣本數 $n$      & 120  & --   \\
    平均值 $\bar{x}$ & 3.14 & cm   \\
    標準差 $s$      & 0.27 & cm   \\
    \bottomrule
  \end{tabular}
\end{table}
